%!TEX root = ../template.tex
%%%%%%%%%%%%%%%%%%%%%%%%%%%%%%%%%%%%%%%%%%%%%%%%%%%%%%%%%%%%%%%%%%%
%% chapter1.tex
%% NOVA thesis document file
%%
%% Chapter with introduction
%%%%%%%%%%%%%%%%%%%%%%%%%%%%%%%%%%%%%%%%%%%%%%%%%%%%%%%%%%%%%%%%%%%

\typeout{NT FILE chapter1.tex}%

\chapter{Introduction}
\label{cha:introduction}

\section{Motivation}
\label{sec:motivation}
Nowadays, 5G networks are no longer just an evolution of mobile broadband communications, they are the technological foundation for a new generation of applications that demand infallible reliability, ultra-low latency and massive connectivity. From industrial automation \cites{13-electronics11030412}{brown2018ultrareliable}{Khujamatov2021}, extended and augmented reality \cites{LV2020158}{9363323}{16-TorresVega2020} to \gls{IoT} applications \cites{7835795}{LV2020158}{16-TorresVega2020}{8519960}, these modern use cases require a robust network architecture that can be reconfigured dynamically, scaled up or down and overall, be tailored to meet their specific needs. 

This shift has placed enormous emphasis on programmability, softwarization and virtualization across both \glsxtrshort{RAN} and Core components. As a result, network operators, research institutions, academic investigators and industry stakeholders are steering away from proprietary systems and looking towards open-source, cloud-native 5G solutions that can be easily adapted to different scenarios and requirements. These platforms make it possible to explore new algorithms, validate architectural choices and test emerging technologies without the hurdles imposed by closed, vendor-specific systems.

However, deploying a fully functional 5G network based solely on open-source components is a complex task, requiring a deep understanding of the 5G architecture, protocols and best practices, as well as familiarity with the plethora of available open-source frameworks, along with their constraints. This report aims to explore the available open-source solutions for 5G networks, comparing their features and drawbacks. 

\subsection{Objectives and Contributions}
\label{sec:objectives_contributions}
This subsection outlines the primary objectives that were defined for this thesis project, as well as a few expected contributions to the field of 5G networking.\\
\underline{Objectives:}
\begin{itemize}
    \item Elaboration of a comprehensive literature review on 5G network architectures, protocols, and open-source implementations, identifying key design challenges and opportunities in the field.
    \item Development of a 5G network testbed, including both \glsxtrshort{RAN} and Core components, using open-source software and accessible hardware.
    \item Evaluation of \glspl{KPI} of the deployed 5G network under various scenarios, such as throughput, latency, reliability and handover performance.
\end{itemize}
\underline{Expected Contributions:}
\begin{itemize}
    \item A detailed analysis of existing open-source 5G solutions, highlighting their strengths, weaknesses, and suitability for different deployment scenarios.
    \item A functional 5G network testbed that can be used for further research and experimentation in the field of 5G networking.
\end{itemize}
\section{Report Structure}
\label{sec:report_structure}
The structure of this document is as follows:
\begin{description}{}
    \item{\uline{Chapter 2}} This chapter details extensively the fundamental elements and architecture of 5G networks, going into detail about both the \gls{RAN} and Core components, along with their features, as well as current open-source solutions for implementing the architecture. An overview of the existing testbeds and experimental projects in the field of 5G networks is also provided. 
    \item{\uline{Chapter 3}} This chapter describes the work plan and the milestones to reach along the thesis project.
    \item{\uline{Chapter 4}} This chapter concludes the document, summarizing the key points and outlining the next steps for the thesis project.
\end{description}
\clearpage